\documentclass[12pt,a4paper]{article}
\usepackage{geometry}
\geometry{left=2.5cm,right=2.5cm,top=2.0cm,bottom=2.5cm}
\usepackage[english]{babel}
\usepackage{amsmath,amsthm}
\usepackage{amsfonts}
\usepackage[longend,ruled,linesnumbered]{algorithm2e}
\usepackage{fancyhdr}
\usepackage{ctex}
\usepackage{array}
\usepackage{listings}
\usepackage{color}
\usepackage{graphicx}

\begin{document}


\title{
{《优化理论与算法》第7次作业
}
}
\date{}

% \author{
% 姓名：\underline{你的名字}~~~~~~
% 学号：\underline{你的学号}~~~~~~}

\maketitle
\begin{itemize}
	\item 作业截止于{\bf 周五（4/24/2025）}，在Canvas系统中提交PDF或照片文件
	\item 作业题目的tex文件可以在Canvas中下载
	\item 鼓励相互讨论，但不要抄袭.
\end{itemize}
%%%%%%%%%%%%%%%%%%%%%%%%%%%%%%%%%%%%%%%%%%%%%%%%%%%%%%%%%%%%%%%%%
%%%%%%%%%%%%%%%%%%%%%%%%%%%%%%%%%%%%%%%%%%%%%%%%%%%%%%%%%%%%%%%%%

\section{阅读与积累}
\paragraph{1.1}
\begin{itemize}
	\item 复习 Boyd and Vandenberghe，也就是 Convex Optimization 教材，章节 2.1-2.5 与 3.1相关内容。预习章节3.2-3.4。
\item 收看同名超星线上慕课章节2.1内容。
\end{itemize}

\section{习题}
\paragraph{2.1}
设$C\subseteq\mathbb{R}^n$为一个凸集且$x_1,...,x_k \in C$。令$\theta_1,...,\theta_k \in \mathbb{R}$满足$\theta_i \geq 0, \;\sum_{i=1}^k \theta_i=1$。证明$\sum_{i=1}^k\theta_i x_i\in C$。（凸性定义是指此式在$k=2$时成立；你需证明对任意$k$的情况。）提示对$k$归纳。


\paragraph{2.2}
下面集合是否为凸集合？如果是，则证明之；如果不是，请找出反例。
\begin{itemize}
\item[(a)] $C=\{x \in \mathbb{R}^n: x^Tx \leq 1\}$
    \item[(b)] $C =\{(x_1,x_2) \in \mathbb{R}^2 : x_1 x_2 \leq 1, x_1 < 0\}$ 
    \item[(c)] $C=\{A \in \mathbb{R}^{4 \times 4} : \text{rank}(A) \leq  2\}$
\end{itemize}

\paragraph{2.3}
下面集合是否为凸集合？如果是，则证明之；如果不是，请找出反例。

\begin{itemize}
  \item[(a)] 平板，即形如 $\{x \in \mathbb{R}^n \mid \alpha \leq a^T x \leq \beta\}$ 的集合。
  
  \item[(b)] 矩形，即形如 $\{x \in \mathbb{R}^n \mid \alpha_i \leq x_i \leq \beta_i,\ i = 1, \dots, n\}$ 的集合。当 $n > 2$ 时，矩形有时也称为超矩形。
  
  \item[(c)] 楔形，即 $\{x \in \mathbb{R}^n \mid a_1^T x \leq b_1,\ a_2^T x \geq b_2\}$。
  
  \item[(d)] 距离给定点比距离给定集合近的点构成的集合，即
  \[
  \{x \mid \|x - x_0\|_2 \leq \|x - y\|_2,\ \forall y \in S\},
  \]
  其中 $S \subset \mathbb{R}^n$。
  
  \item[(e)] 距离一个集合比另一个集合更近的点的集合，即
  \[
  \{x \mid \mathrm{dist}(x, S) \leq \mathrm{dist}(x, T)\},
  \]
  其中 $S, T \subset \mathbb{R}^n$，\\
  $\mathrm{dist}(x, S) = \inf\{\|x - z\|_2 \mid z \in S\}$。
\end{itemize}



\paragraph{2.4}
下面函数是否为凸函数？如果是，则证明之；如果不是，请找出反例。
\begin{enumerate}
	\item $f(x_1,x_2)=x_1 x_2$ on $\mathbb{R}_{++}^2$
	
	\item $f(x_1,x_2)=1/(x_1x_2)$  on $\mathbb{R}_{++}^2$
	
	\item $f(x_1,x_2)=-x_1^\alpha x_2^{1-\alpha}$  on $\mathbb{R}_{++}^2$, where $0 \leq \alpha \leq 1$
\end{enumerate}

\paragraph{2.5}
正式证明《凸函数》课件P19页，函数$f(x)$为凸函数与函数$g(t)$为凸函数的等价性。

\end{document}